\subsection{Positive Operators}

\begin{exercise}{7c1}
    Suppose $T \in \lin{V}$.
    Prove that if both $T$ and $-T$ are positive operators, then $T = 0$.
\end{exercise}

\begin{sol}
    For every $v \in V$, we have $\gen{Tv, v} \geq 0$ and $\gen{-Tv, v} = -\gen{Tv, v} \geq 0$.
    Both inequalities force $\gen{Tv, v} = 0$ for every $v \in V$.
    By 7.43, $Tv = 0$ for every $v \in V$, so $T = 0$.
\end{sol}

\begin{exercise}{7c2}
    Suppose $T \in \lin{\fbb^4}$ is the operator whose matrix (with respect to the standard basis) is
    \[
        \begin{pmatrix}
            2 & -1 & 0 & 0 \\
            -1 & 2 & -1 & 0 \\
            0 & -1 & 2 & -1 \\
            0 & 0 & -1 & 2
        \end{pmatrix}.
    \]
    Show that $T$ is an invertible positive operator.
\end{exercise}

\begin{sol}
    Note that the standard basis is orthonormal, and $\mcal(T)$ as given is equal to its conjugate transpose, so $T$ is self-adjoint.
    For any $x = (x_1, x_2, x_3, x_4) \in \fbb^4$, we have
    \[
        T x = (2x_1 - x_2, -x_1 + 2x_2 - x_3, -x_2 + 2x_3 - x_4, -x_3 + 2x_4).
    \]
    Since $x_i \bar{x_j} + \bar{x_i} x_j = 2 \Re(x_i \bar{x_j})$, we have
    \begin{align*}
        \gen{T x, x} &= (2x_1 - x_2)\bar{x_1} + (-x_1 + 2x_2 - x_3)\bar{x_2} + (-x_2 + 2x_3 - x_4)\bar{x_3} + (-x_3 + 2x_4)\bar{x_4} \\
        &=2|x_1|^2 - \bar{x_1} x_2 - x_1 \bar{x_2} + 2|x_2|^2 - \bar{x_2} x_3 - x_2 \bar{x_3} + 2|x_3|^2 - \bar{x_3} x_4 - x_3 \bar{x_4} + 2|x_4|^2 \\
        &=2|x_1|^2 - 2\Re(\bar{x_1} x_2) + 2|x_2|^2 - 2\Re(\bar{x_2} x_3) + 2|x_3|^2 - 2\Re(\bar{x_3} x_4) + 2|x_4|^2 \\
        &=|x_1|^2 + |x_1 - x_2|^2 + |x_2 - x_3|^2 + |x_3 - x_4|^2 + |x_4|^2 \geq 0.
    \end{align*}
    Hence, $T$ is a positive operator.
    Moreover, if $T x = 0$ for some $x \in \fbb^4$, then
    \[
        0 = \gen{T x, x} = |x_1|^2 + |x_1 - x_2|^2 + |x_2 - x_3|^2 + |x_3 - x_4|^2 + |x_4|^2.
    \]
    This forces $x_1 = x_4 = 0$ so that $x_2 = x_3 = 0$ as well.
    Hence, $x = 0$, and therefore $T$ is invertible by 3.65.
\end{sol}

\begin{exercise}{7c3}
    Suppose $n$ is a positive integer and $T \in \lin{\fbb^n}$ is the operator whose matrix (with respect to the standard basis) consists of all $1$'s.
    Show that $T$ is a positive operator.
\end{exercise}

\begin{sol}
    Let $x = (1, \ldots, 1) \in \fbb^n$.
    For any $y = \tuple y \in \fbb^n$, we have
    \[
        T y = \left(\sum_{i=1}^n y_i, \ldots, \sum_{i=1}^n y_i\right) = \gen{y, x} x.
    \]
    Consider the operator $R \in \lin{\fbb^n}$ defined by $R z = (\gen{z, x}, 0, \ldots, 0)$ for any $z \in \fbb^n$.
    Then
    \[
        \gen{R z, y} = \gen{z, x} \bar{y_1} = \gen{z, y_1 x}
    \]
    so that $R^\ast y = y_1 x$.
    Therefore, $R^\ast R y = \gen{y, x} x = T y$ for any $y \in \fbb^n$.
    Hence, $T = R^\ast R$, which shows that $T$ is a positive operator by 7.38(f).
\end{sol}

\begin{exercise}{7c4}
    Suppose $n$ is an integer with $n > 1$.
    Show that there exists an $n$-by-$n$ matrix $A$ such that all of the entries of $A$ are positive numbers and $A = A^\ast$, but the operator on $\fbb^n$ whose matrix (with respect to the standard basis) equals $A$ is not a positive operator.
\end{exercise}

\begin{sol}
    Let $A$ be the $n$-by-$n$ matrix whose entries are all $1$ except $A_{1, 2} = A_{2, 1} = 2$.
    Then all entries of $A$ are positive numbers, and $A = A^\ast$ by construction.
    Let $T \in \lin{\fbb^n}$ be the operator whose matrix with respect to the standard basis is $A$.
    However, consider $x = (1, -1, 0, \ldots, 0) \in \fbb^n$.
    Then $(T x)_1 = 1 + 2(-1) = -1$, $(T x)_2 = 2(1) + (-1) = 1$, and $(T x)_i = 1 - 1 = 0$ for $i > 2$.
    Hence,
    \[
        \gen{T x, x} = (-1)(1) + (1)(-1) + 0 = -2
    \]
    so that $T$ is not a positive operator.
\end{sol}

\begin{exercise}{7c5}
    Suppose $T \in \lin{V}$ is self-adjoint.
    Prove that $T$ is a positive operator if and only if for every orthonormal basis $e_1, \ldots, e_n$ of $V$, all entries on the diagonal of $\mcal(T, (e_1, \ldots, e_n))$ are nonnegative numbers.
\end{exercise}

\begin{sol}
    Note that the $(j, k)$th entry of $\mcal(T, (e_1, \ldots, e_n))$ is $\gen{T e_k, e_j}$, hence the diagonal entries are $\gen{T e_j, e_j}$ for $j = 1, \ldots, n$.

    \rightimp Since $T$ is self-adjoint by hypothesis, $T$ is positive if and only if $\gen{T v, v} \geq 0$ for all $v \in V$.
    In particular, this must hold for each $v = e_j$ in any orthonormal basis $e_1, \ldots, e_n$ of $V$, which shows that all diagonal entries of $\mcal(T, (e_1, \ldots, e_n))$ are nonnegative.

    \leftimp If $v = 0$, then $\gen{T v, v} = 0 \geq 0$.
    Suppose $v \neq 0$, and let $e_1 = v / \norm{v}$.
    Then $\norm{e_1} = 1$, so $e_1$ is an orthonormal list by itself.
    Extend $e_1$ to an orthonormal basis $e_1, \ldots, e_n$ of $V$.
    Then the first diagonal entry of $\mcal(T, (e_1, \ldots, e_n))$ is $\gen{T e_1, e_1} = \gen{T (v / \norm{v}), v / \norm{v}} = \gen{T v, v} / \norm{v}^2 \geq 0$.
    Hence, $\gen{T v, v} \geq 0$ for all $v \in V$, which shows that $T$ is a positive operator.
\end{sol}

\begin{exercise}{7c6}
    Prove that the sum of two positive operators on $V$ is a positive operator.
\end{exercise}

\begin{sol}
    Let $S, T \in \lin{V}$ be positive operators.
    In particular, $S$ and $T$ are self-adjoint.
    By 7.5(a), $S + T$ is also self-adjoint.
    Positivity follows from the fact that for any $v \in V$,
    \[
        \gen{(S + T)v, v} = \gen{Sv, v} + \gen{Tv, v} \geq 0 + 0 = 0,
    \]
    since $S$ and $T$ are positive operators.
    Hence, $S + T$ is a positive operator.
\end{sol}

\begin{exercise}{7c7}
    Suppose $S \in \lin{V}$ is an invertible positive operator and $T \in \lin{V}$ is a positive operator.
    Prove that $S + T$ is invertible.
\end{exercise}

\begin{sol}
    If $(S + T)v = 0$ for some $v \in V$, then $\gen{Sv, v} + \gen{Tv, v} = \gen{(S + T)v, v} = \gen{0, v} = 0$.
    Since $S$ and $T$ are positive, $\gen{Sv, v} \geq 0$ and $\gen{Tv, v} \geq 0$.
    In particular, $\gen{Sv, v} = 0$.
    By 7.43, $Sv = 0$ so that $v = 0$ by the invertibility of $S$.
    Therefore, $(S + T)v = 0$ only if $v = 0$, which shows that $S + T$ is invertible by 3.65.
\end{sol}

\begin{exercise}{7c8}
    Suppose $T \in \lin{V}$.
    Prove that $T$ is a positive operator if and only if the pseudoinverse $T^\dagger$ is a positive operator.
\end{exercise}

\begin{soliff}
    \rightimp By 7.38(c), there is an orthonormal basis $e_1, \ldots, e_n$ of $V$ consisting of eigenvectors of $T$ with respect to which the matrix of $T$ is diagonal with nonnegative diagonal entries $\lambda_1, \ldots, \lambda_n$.
    Hence, $T e_j = \lambda_j e_j$, where $\lambda_j \geq 0$ for each $j$.
    By \exref{7b25},
    \[
        T^\dagger e_j = 
        \begin{cases}
            \frac{1}{\lambda_j} e_j & \text{if } \lambda_j \neq 0, \\
            0 & \text{if } \lambda_j = 0.
        \end{cases}
    \]
    Hence, $\mcal(T^\dagger, (e_1, \ldots, e_n))$ is a diagonal matrix with nonnegative entries, which shows that $T^\dagger$ is a positive operator by 7.38(c).

    \leftimp By \exref{6c23}, $(T^\dagger)^\dagger = T$. 
    Applying the forward direction to $T^\dagger$ in place of $T$ shows that $T$ is positive.
\end{soliff}

\begin{exercise}{7c9}
    Suppose $T \in \lin{V}$ is a positive operator and $S \in \lin{W, V}$.
    Prove that $S^\ast T S$ is a positive operator on $W$.
\end{exercise}

\begin{sol}
    Note that $S^\ast \in \lin{V, W}$, so $S^\ast T S \in \lin{W}$.
    By 7.5(c) and 7.5(d), we have $(S^\ast T S)^\ast = S^\ast T^\ast (S^\ast)^\ast = S^\ast T S$, where the last equality holds since $T$ is self-adjoint.
    Moreover, for any $w \in W$,
    \[
        \gen{S^\ast T S w, w} = \gen{T S w, S w} \geq 0,
    \]
    since $T$ is positive.
    Hence, $S^\ast T S$ is a positive operator on $W$.
\end{sol}

\begin{exercise}{7c10}
    Suppose $T$ is a positive operator on $V$.
    Suppose $v, w \in V$ are such that
    \[
        Tv = w \quad \text{and} \quad Tw = v.
    \]
    Prove that $v = w$.
\end{exercise}

\begin{sol}
    Consider $v - w$.
    Then
    \[
        \gen{T(v - w), v - w} = \gen{Tv - Tw, v - w} = \gen{w - v, v - w} = -\gen{v - w, v - w} = -\norm{v - w}^2.
    \]
    Since $T$ is positive, $\gen{T(v - w), v - w} \geq 0$, so $-\norm{v - w}^2 \geq 0$, which implies $\norm{v - w}^2 = 0$ and hence $v - w = 0$, i.e., $v = w$.
\end{sol}

\begin{exercise}{7c11}
    Suppose $T$ is a positive operator on $V$ and $U$ is a subspace of $V$ invariant under $T$.
    Prove that $T|_U \in \lin{U}$ is a positive operator on $U$.
\end{exercise}

\begin{sol}
    Note that $T|_U \in \lin{U}$ since $U$ is invariant under $T$, and $U$ carries the inner product inherited from $V$.
    Since $T$ is self-adjoint on $V$, \exref{7b19}(b) gives that $T|_U$ is self-adjoint on $U$.
    Moreover, for any $u \in U$,
    \[
        \gen{T|_U u, u}_U = \gen{Tu, u}_V \geq 0,
    \]
    since $T$ is positive on $V$.
    Hence, $T|_U$ is a positive operator on $U$.
\end{sol}

\begin{exercise}{7c12}
    Suppose $T \in \lin{V}$ is a positive operator.
    Prove that $T^k$ is a positive operator for every positive integer $k$.
\end{exercise}

\begin{sol}
    By 7.38(c), there is an orthonormal basis $e_1, \ldots, e_n$ of $V$ with respect to which $\mcal(T, (e_1, \ldots, e_n))$ is diagonal with nonnegative entries $\lambda_1, \ldots, \lambda_n$.
    By 3.43 and induction on $k$, $\mcal(T^k, (e_1, \ldots, e_n)) = (\mcal(T, (e_1, \ldots, e_n)))^k$ is diagonal with entries $\lambda_1^k, \ldots, \lambda_n^k$, which are all nonnegative.
    By 7.38(c) again, this shows that $T^k$ is positive.
\end{sol}

\begin{exercise}{7c13}
    Suppose $T \in \lin{V}$ is self-adjoint and $\alpha \in \rbb$.
    \begin{subexercises}
        \item Prove that $T - \alpha I$ is a positive operator if and only if $\alpha$ is less than or equal to every eigenvalue of $T$.
        \item Prove that $\alpha I - T$ is a positive operator if and only if $\alpha$ is greater than or equal to every eigenvalue of $T$.
    \end{subexercises}
\end{exercise}

\begin{solenum}
    \item \rightimp If $T - \alpha I$ is positive, then its eigenvalues are nonnegative by 7.38(b).
    Since $Tv = \lambda v$ if and only if $(T - \alpha I)v = (\lambda - \alpha)v$, the eigenvalues of $T - \alpha I$ are precisely $\lambda - \alpha$, where $\lambda$ ranges over the eigenvalues of $T$.
    Hence, $\alpha \leq \lambda$ for every eigenvalue $\lambda$ of $T$.

    \leftimp By 7.5(a), (b), (e), self-adjointness of $T$, and $\alpha \in \rbb$,
    \[
        (T - \alpha I)^\ast = T^\ast - \bar\alpha I^\ast = T - \alpha I,
    \]
    so $T - \alpha I$ is self-adjoint.
    Since $\lambda - \alpha \geq 0$ for every eigenvalue $\lambda$ of $T$, the eigenvalues of $T - \alpha I$ are nonnegative.
    Hence, $T - \alpha I$ is positive by 7.38(b).
    \item Note that $-T$ is self-adjoint by 7.5(b).
    Moreover, the eigenvalues of $-T$ are exactly $-\lambda$, where $\lambda$ ranges over the eigenvalues of $T$.
    Then applying part (a) to $-T$ and $-\alpha$ gives the operator $-T - (-\alpha)I = \alpha I - T$, which is positive if and only if $-\alpha \leq -\lambda$ for every eigenvalue $\lambda$ of $T$, i.e., $\alpha \geq \lambda$ for every eigenvalue $\lambda$ of $T$.
\end{solenum}

\begin{exercise}{7c14}
    Suppose $T$ is a positive operator on $V$ and $v_1, \ldots, v_m \in V$.
    Prove that
    \[
        \sum_{j=1}^{m} \sum_{k=1}^{m} \gen{Tv_k, v_j} \geq 0.
    \]
\end{exercise}

\begin{sol}
    Let $v = \sum_{k=1}^{m} v_k$.
    By linearity of $T$, additivity of the inner product, and positivity of $T$, we have
    \[
        \sum_{j=1}^{m} \sum_{k=1}^{m} \gen{Tv_k, v_j} = \sum_{j=1}^{m} \gen{T v, v_j} = \gen{Tv, v} \geq 0. \qedhere
    \]
\end{sol}

\begin{exercise}{7c15}
    Suppose $T \in \lin{V}$ is self-adjoint.
    Prove that there exist positive operators $A, B \in \lin{V}$ such that
    \[
        T = A - B \quad \text{and} \quad \sqrt{T^\ast T} = A + B \quad \text{and} \quad AB = BA = 0.
    \]
\end{exercise}

\begin{sol}
    If $\fbb = \rbb$, the real spectral theorem applies directly, and if $\fbb = \cbb$, self-adjointness implies normality, so the complex spectral theorem applies.
    In either case, there exists an orthonormal basis $e_1, \ldots, e_n$ of $V$ consisting of eigenvectors of $T$ with corresponding eigenvalues $\lambda_1, \ldots, \lambda_n$.
    Let $\lambda_k^+ = \max\{\lambda_k, 0\}$ and $\lambda_k^- = \max\{-\lambda_k, 0\}$ for each $k = 1, \ldots, n$.
    By 3.4, there exist operators $A$ and $B$ on $V$ such that
    \[
        A e_k = \lambda_k^+ e_k \quad \text{and} \quad B e_k = \lambda_k^- e_k
    \]
    for each $k = 1, \ldots, n$.
    Since $\mcal(A, (e_1, \ldots, e_n))$ and $\mcal(B, (e_1, \ldots, e_n))$ are diagonal with nonnegative entries, both $A$ and $B$ are positive operators by 7.38(c).
    Since $\lambda_k = \lambda_k^+ - \lambda_k^-$ for each $k$, $T e_k = (\lambda_k^+ - \lambda_k^-) e_k = (A - B) e_k$ for each $k$, so $T = A - B$ by 3.4.
    Additionally, $\lambda_k^+ \lambda_k^- = 0$ for each $k$, so $AB e_k = \lambda_k^+ \lambda_k^- e_k = 0$ for each $k$, and similarly $BA e_k = 0$ for each $k$, hence $AB = BA = 0$ by 3.4.
    
    Finally, note that since $T$ is self-adjoint, $T^\ast T = T^2$, so $\sqrt{T^\ast T} = \sqrt{T^2}$.
    By 3.4, there exists an operator $C$ on $V$ such that $C e_k = |\lambda_k| e_k$ for each $k = 1, \ldots, n$.
    Since $\mcal(C, (e_1, \ldots, e_n))$ is diagonal with nonnegative entries, $C$ is a positive operator by 7.38(c).
    Moreover, $C^2 e_k = |\lambda_k|^2 e_k = \lambda_k^2 e_k = T^2 e_k$ for each $k$, so $C^2 = T^2$ by 3.4.
    By 7.39, $C = \sqrt{T^2} = \sqrt{T^\ast T}$.
    Also, note that $C e_k = |\lambda_k| e_k = (\lambda_k^+ + \lambda_k^-) e_k = (A + B) e_k$ for each $k$, so $C = A + B$ by 3.4.
\end{sol}

\begin{exercise}{7c16}
    Suppose $T$ is a positive operator on $V$.
    Prove that
    \[
        \nul \sqrt{T} = \nul T \quad \text{and} \quad \range \sqrt{T} = \range T.
    \]
\end{exercise}

\begin{sol}
    Since $T$ is positive, $\sqrt T$ is also a positive operator.
    In particular, it is self-adjoint, hence normal.
    Let $S = \sqrt{T}$, so $S^2 = T$.
    By \exref{7a27}, $\nul S = \nul S^2 = \nul T$ and $\range S = \range S^2 = \range T$.
\end{sol}

\begin{exercise}{7c17}
    Suppose that $T \in \lin{V}$ is a positive operator.
    Prove that there exists a polynomial $p$ with real coefficients such that $\sqrt{T} = p(T)$.
\end{exercise}

\begin{sol}
    If $V = \set 0$, take $p = 0$.
    Otherwise, $V \neq \set 0$.
    By 7.38(c), there exists an orthonormal basis $e_1, \ldots, e_n$ of $V$ such that $T e_k = \lambda_k e_k$ where $\lambda_k \geq 0$ for each $k = 1, \ldots, n$.
    Since $n \geq 1$, $T$ has at least one eigenvalue.
    Let $\mu_1, \ldots, \mu_m$ be the distinct eigenvalues of $T$, so $m \geq 1$.
    Because $e_k \neq 0$, $\lambda_k$ is an eigenvalue of $T$ so that $\lambda_k \in \set{\mu_1, \ldots, \mu_m}$ for each $k = 1, \ldots, n$.
    By 7.38(b), each $\mu_j \geq 0$ so that $\sqrt{\mu_j} \in \rbb$ for each $j = 1, \ldots, m$.
    By \exref{4-7} applied over $\rbb$, there exists $p \in \pcal_{m-1}(\rbb)$ such that $p(\mu_j) = \sqrt{\mu_j}$ for each $j = 1, \ldots, m$.
    Since $T^i e_k = \lambda_k^i e_k$ for each nonnegative integer $i$, the definition of $p(T)$ gives $p(T) e_k = p(\lambda_k) e_k$ for each $k = 1, \ldots, n$.
    Since $\lambda_k \in \set{\mu_1, \ldots, \mu_m}$ and $p(\mu_j) = \sqrt{\mu_j}$ for each $j = 1, \ldots, m$, we have $p(\lambda_k) = \sqrt{\lambda_k}$ for each $k = 1, \ldots, n$.
    Then $\mcal(p(T), (e_1, \ldots, e_n))$ is diagonal with diagonal entries $\sqrt{\lambda_1}, \ldots, \sqrt{\lambda_n}$.
    By 7.38(c), $p(T)$ is positive.
    Then
    \[
        p(T)^2 e_k = p(\lambda_k)^2 e_k = (\sqrt{\lambda_k})^2 e_k = \lambda_k e_k = T e_k
    \]
    for each $k = 1, \ldots, n$.
    By 3.4, $p(T)^2 = T$, and by 7.39, $p(T) = \sqrt{T}$.
\end{sol}

\begin{exercise}{7c18}
    Suppose $S$ and $T$ are positive operators on $V$.
    Prove that $ST$ is a positive operator if and only if $S$ and $T$ commute.
\end{exercise}

\begin{soliff}
    \rightimp Note that $S$, $T$, and $ST$ are all self-adjoint.
    By \exref{7a9}, $ST$ is self-adjoint if and only if $ST = TS$.

    \leftimp By \exref{7c17}, there exists a polynomial $p$ such that $\sqrt S = p(S)$.
    Since $T$ commutes with $S$, it commutes with powers of $S$ so that $T$ commutes with $p(S) = \sqrt{S}$.
    Then $ST = \sqrt{S} \sqrt{S} T = \sqrt{S}^\ast \sqrt{S} T = \sqrt{S}^\ast T \sqrt{S}$, where the second equality uses the fact that $\sqrt{S}$ is self-adjoint.
    By \exref{7c9}, $ST$ is positive.
\end{soliff}

\begin{exercise}{7c19}
    Show that the identity operator on $\fbb^2$ has infinitely many self-adjoint square roots.
\end{exercise}

\begin{sol}
    Let $x, y \in \rbb$, and define the family of operators $T_{x, y}$ on $\fbb^2$ whose matrices with respect to the standard basis are
    \[
        \mcal(T_{x, y}) =
        \begin{pmatrix}
            x & y \\
            y & -x
        \end{pmatrix}.
    \]
    Then $\mcal(T_{x, y}) = \mcal(T_{x, y})^*$.
    Since the standard basis is orthonormal, $T_{x, y}$ is self-adjoint by 7.9.
    Moreover,
    \[
        \mcal(T_{x, y})^2 =
        \begin{pmatrix}
            x & y \\
            y & -x
        \end{pmatrix}^2
        =
        \begin{pmatrix}
            x^2 + y^2 & 0 \\
            0 & x^2 + y^2
        \end{pmatrix} = (x^2 + y^2) I.
    \]
    Taking $x = \cos \theta$ and $y = \sin \theta$ for $\theta \in \rbb$, we get $x^2 + y^2 = 1$.
    Since $(\cos \theta, \sin \theta)$ takes on infinitely many distinct values and $T_{\cos \theta, \sin \theta}(1, 0) = (\cos \theta, \sin \theta)$, the operators $T_{\cos \theta, \sin \theta}$ are all distinct.
    Therefore, the identity operator on $\fbb^2$ has infinitely many self-adjoint square roots.
\end{sol}

\begin{exercise}{7c20}
    Suppose $T \in \lin{V}$ and $e_1, \ldots, e_n$ is an orthonormal basis of $V$.
    Prove that $T$ is a positive operator if and only if there exist $v_1, \ldots, v_n \in V$ such that
    \[
        \gen{Te_k, e_j} = \gen{v_k, v_j}
    \]
    for all $j, k = 1, \ldots, n$.

    \note{The numbers $\set{\gen{Te_k, e_j}}_{j,k=1,\ldots,n}$ are the entries in the matrix of $T$ with respect to the orthonormal basis $e_1, \ldots, e_n$.}
\end{exercise}

\begin{soliff}
    \rightimp By 7.38(f), let $T = R^\ast R$ for some $R \in \lin{V}$. 
    Define $v_k = R e_k$ for $k = 1, \ldots, n$. 
    Then for all $j, k = 1, \ldots, n$, we have
    \[
        \gen{v_k, v_j} = \gen{R e_k, R e_j} = \gen{R^\ast R e_k, e_j} = \gen{T e_k, e_j}.
    \]
    \leftimp By 3.4, we may define the unique operator $R \in \lin{V}$ such that $R e_k = v_k$ for $k = 1, \ldots, n$.
    Then
    \[
        \gen{T e_k, e_j} = \gen{v_k, v_j} = \gen{R e_k, R e_j} = \gen{R^\ast R e_k, e_j}.
    \]
    Let $S = T - R^\ast R$.
    Then
    \[
        \gen{S e_k, e_j} = 0
    \]
    for all $j, k = 1, \ldots, n$.
    By 6.30(a), we have
    \[
        S e_k = \sum_{j=1}^n \gen{S e_k, e_j} e_j = 0.
    \]
    Since this holds for all $k = 1, \ldots, n$, we have $S = 0$ by 3.4 so that $T = R^\ast R$.
    Therefore, $T$ is positive by 7.38(f).
\end{soliff}

\begin{exercise}{7c21}
    Suppose $n$ is a positive integer.
    The $n$-by-$n$ \term{Hilbert matrix} is the $n$-by-$n$ matrix whose entry in row $j$, column $k$ is $\frac{1}{j+k-1}$.
    Suppose $T \in \lin{V}$ is an operator whose matrix with respect to some orthonormal basis of $V$ is the $n$-by-$n$ Hilbert matrix.
    Prove that $T$ is a positive invertible operator.

    \note{Example: The $4$-by-$4$ Hilbert matrix is
    \[
        \begin{pmatrix}
            1 & \frac{1}{2} & \frac{1}{3} & \frac{1}{4} \\
            \frac{1}{2} & \frac{1}{3} & \frac{1}{4} & \frac{1}{5} \\
            \frac{1}{3} & \frac{1}{4} & \frac{1}{5} & \frac{1}{6} \\
            \frac{1}{4} & \frac{1}{5} & \frac{1}{6} & \frac{1}{7}
        \end{pmatrix}.
    \]}
\end{exercise}

\begin{sol}
    Let $e_1, \ldots, e_n$ be the orthonormal basis of $V$ with respect to which the matrix of $T$ is the $n$-by-$n$ Hilbert matrix.
    Let $H_n = \mcal(T, (e_1, \ldots, e_n))$.
    Since $H_n$ has real entries, and $\frac{1}{j+k-1} = \frac{1}{k+j-1}$, it follows that $H_n = H_n^\ast$.
    By 7.9, it follows that $T$ is self-adjoint.

    Now, observe that
    \[
        \int_0^1 x^{j-1} x^{k-1} \dd x = \frac{1}{j+k-1}
    \]
    for all $j, k = 1, \ldots, n$.
    Consider $v = \sum_{k=1}^n \alpha_k e_k \in V$ for $\alpha_1, \ldots, \alpha_n \in \fbb$, and define the polynomial $p(x) = \sum_{k=1}^n \alpha_k x^{k-1}$.
    Note that $\bar{p(x)} = \sum_{k=1}^n \bar{\alpha_k} x^{k-1}$ since $x$ is real.
    Then
    \begin{align*}
        \gen{T v, v} &= \sum_{j=1}^n \sum_{k=1}^n \alpha_k \bar{\alpha_j} \gen{T e_k, e_j} \\
        &= \sum_{j=1}^n \sum_{k=1}^n \alpha_k \bar{\alpha_j} \frac{1}{j+k-1} \\
        &= \sum_{j=1}^n \sum_{k=1}^n \alpha_k \bar{\alpha_j} \int_0^1 x^{j-1} x^{k-1} \dd x \\
        &= \int_0^1 \sum_{j=1}^n \sum_{k=1}^n \alpha_k \bar{\alpha_j} x^{j-1} x^{k-1} \dd x \\
        &= \int_0^1 \sum_{k=1}^n \alpha_k x^{k-1} \sum_{j=1}^n \bar{\alpha_j} x^{j-1} \dd x \\
        &= \int_0^1 p(x) \bar{p(x)} \dd x = \int_0^1 \abs{p(x)}^2 \dd x \geq 0.
    \end{align*}
    Hence, $T$ is positive.
    Moreover, suppose $T v = 0$.
    Then $\gen{T v, v} = 0$, which implies $\int_0^1 \abs{p(x)}^2 \dd x = 0$.
    Note that $\abs{p(x)}^2$ is a continuous and nonnegative function on $[0, 1]$, so $\int_0^1 \abs{p(x)}^2 \dd x = 0$ implies that $\abs{p(x)}^2 = 0$ for all $x \in [0, 1]$.
    Then $p$ has infinitely many zeros in $[0, 1]$ while a nonzero polynomial of degree at most $n-1$ can have at most $n-1$ zeros.
    Therefore, $p$ is the zero polynomial.
    Hence, $\alpha_1 = \cdots = \alpha_n = 0$, so $v = 0$.
    Since $V$ is finite-dimensional, it follows that $T$ is invertible by 3.65.
\end{sol}

\begin{exercise}{7c22}
    Suppose $T \in \lin{V}$ is a positive operator and $u \in V$ is such that $\norm{u} = 1$ and $\norm{Tu} \geq \norm{Tv}$ for all $v \in V$ with $\norm{v} = 1$.
    Show that $u$ is an eigenvector of $T$ corresponding to the largest eigenvalue of $T$.
\end{exercise}

\begin{sol}
    Since $u \in V$, $V \neq \set 0$.
    By 7.38(c), there exists an orthonormal basis $e_1, \ldots, e_n$ of $V$ consisting of eigenvectors of $T$ with nonnegative eigenvalues $\lambda_1, \ldots, \lambda_n$.
    Define
    \[
        \lambda = \max_{1 \leq k \leq n} \lambda_k.
    \]
    For any $v \in V$ with $\norm{v} = 1$, write $v = \sum_{k=1}^n \alpha_k e_k$.
    Then
    \[
        \norm{T v}^2 = \norm*{\sum_{k=1}^n \alpha_k \lambda_k e_k}^2 = \sum_{k=1}^n \abs{\alpha_k}^2 \lambda_k^2 \leq \lambda^2 \sum_{k=1}^n \abs{\alpha_k}^2 = \lambda^2 \norm{v}^2 = \lambda^2.
    \]
    Since $\norm{Tu} \geq \norm{Tv}$ for all $v \in V$ with $\norm{v} = 1$, the above shows that $\norm{Tu} \leq \lambda$.
    Moreover, if $e_k$ is an eigenvector corresponding to the eigenvalue $\lambda$, then $\norm{T e_k} = \lambda$.
    By hypothesis, $\norm{Tu} \geq \norm{T e_k} = \lambda$, so $\norm{Tu} = \lambda$.
    Let $u = \sum_{k=1}^n \beta_k e_k$.
    Then $\sum_{k=1}^n \abs{\beta_k}^2 = 1$, so
    \[
        \norm{Tu}^2 = \sum_{k=1}^n \abs{\beta_k}^2 \lambda_k^2 = \lambda^2, \quad \text{and} \quad \lambda^2 \sum_{k=1}^n \abs{\beta_k}^2 = \lambda^2.
    \]
    Observe that
    \[
        \sum_{k=1}^n \abs{\beta_k}^2 (\lambda^2 - \lambda_k^2) = 0.
    \]
    Since $\lambda_k \leq \lambda$ for all $k$, equality holds if and only if $\abs{\beta_k}^2 = 0$ for all $k$ with $\lambda_k < \lambda$.
    Hence, $u$ is a linear combination of eigenvectors corresponding to the eigenvalue $\lambda$.
    Therefore, $u$ is an eigenvector of $T$ corresponding to the largest eigenvalue $\lambda$.
\end{sol}

\begin{exercise}{7c23}
    For $T \in \lin{V}$ and $u, v \in V$, define $\gen{u, v}_T$ by $\gen{u, v}_T = \gen{Tu, v}$.
    \begin{subexercises}
        \item Suppose $T \in \lin{V}$.
        Prove that $\gen{\cdot, \cdot}_T$ is an inner product on $V$ if and only if $T$ is an invertible positive operator (with respect to the original inner product $\gen{\cdot, \cdot}$).
        \item Prove that every inner product on $V$ is of the form $\gen{\cdot, \cdot}_T$ for some positive invertible operator $T \in \lin{V}$.
    \end{subexercises}
\end{exercise}

\begin{solenum}
    \item \rightimp Suppose $\gen{\cdot, \cdot}_T$ is an inner product on $V$.
    By positivity, $\gen{Tv, v} \geq 0$ for all $v \in V$.
    By conjugate symmetry,
    \[
        \gen{Tu, v} = \gen{u, v}_T = \bar{\gen{v, u}_T} = \bar{\gen{Tv, u}} = \gen{u, Tv}
    \]
    for all $u, v \in V$, so $T$ is self-adjoint.
    Therefore, $T$ is positive.
    If $T v = 0$ for some $v \in V$, then $\gen{Tv, v} = 0$, which implies $v = 0$ by definiteness of the inner product $\gen{\cdot, \cdot}_T$.
    Hence, $T$ is invertible by 3.65.

    \leftimp Suppose $T$ is an invertible positive operator, and define $\gen{\cdot, \cdot}_T$ as given.
    By 7.38(f), there exists $R \in \lin{V}$ such that $R^\ast R = T$, so
    \[
        \gen{v, v}_T = \gen{Tv, v} = \gen{R^\ast R v, v} = \gen{Rv, Rv} = \norm{Rv}^2
    \]
    for any $v \in V$.
    We now verify the inner product axioms.
    Let $u, v, w \in V$ and $\lambda \in \fbb$.
    Then:
    \begin{itemize}
        \item \textbf{Positivity}: $\gen{v, v}_T = \gen{Tv, v} \geq 0$ since $T$ is positive.
        \item \textbf{Definiteness}: If $\gen{v, v}_T = 0$, then $\norm{Rv}^2 = 0$, which implies $Rv = 0$ and hence $Tv = R^\ast(Rv) = 0$.
        Since $T$ is invertible, this implies $v = 0$.
        Conversely, $\gen{0, 0}_T = \gen{T0, 0} = \gen{0, 0} = 0$.
        \item \textbf{Additivity in the First Slot}: 
        \[
            \gen{u + w, v}_T = \gen{T(u + w), v} = \gen{Tu + Tw, v} = \gen{Tu, v} + \gen{Tw, v} = \gen{u, v}_T + \gen{w, v}_T.
        \]
        \item \textbf{Homogeneity in the First Slot}: 
        \[
            \gen{\lambda u, v}_T = \gen{T(\lambda u), v} = \gen{\lambda Tu, v} = \lambda \gen{Tu, v} = \lambda \gen{u, v}_T.
        \]
        \item \textbf{Conjugate Symmetry}: Since $T$ is self-adjoint, we have
        \[
            \gen{u, v}_T = \gen{Tu, v} = \gen{u, Tv} = \bar{\gen{Tv, u}} = \bar{\gen{v, u}_T}.
        \]
    \end{itemize}
    Hence, $\gen{\cdot, \cdot}_T$ satisfies all the axioms of an inner product on $V$, so it is indeed an inner product.
    \item Let $\gen{\cdot, \cdot}'$ be an inner product on $V$.
    For each $v \in V$, define $\phi_v : V \to \fbb$ by $\phi_v(u) = \gen{u, v}'$.
    Since $\gen{\cdot, \cdot}'$ is an inner product on $V$, $\phi_v$ is a linear functional on $V$.
    By the Riesz representation theorem, there exists a unique $Sv \in V$ such that $\phi_v(u) = \gen{u, Sv}$ for all $u \in V$.
    This defines a function $S : V \to V$.
    We claim that $S$ is linear.
    To see this, let $v, w \in V$ and $\mu \in \fbb$. Then for all $u \in V$,
    \[
        \gen{u, S(v + w)} = \gen{u, v + w}' = \gen{u, v}' + \gen{u, w}' = \gen{u, Sv} + \gen{u, Sw} = \gen{u, Sv + Sw},
    \]
    which implies $S(v + w) = Sv + Sw$ by the uniqueness part of the Riesz representation theorem.
    Similarly,
    \[
        \gen{u, S(\mu v)} = \gen{u, \mu v}' = \bar\mu \gen{u, v}' = \bar\mu \gen{u, Sv} = \gen{u, \mu Sv},
    \]
    which implies $S(\mu v) = \mu Sv$ by the uniqueness part of the Riesz representation theorem.
    Hence, $S$ is linear.
    By conjugate symmetry of both $\gen{\cdot, \cdot}'$ and $\gen{\cdot, \cdot}$, 
    \[
        \gen{u, v}' = \bar{\gen{v, u}'} = \bar{\gen{v, Su}} = \gen{Su, v} = \gen{u, v}_S
    \]
    for all $u, v \in V$, so $\gen{\cdot, \cdot}' = \gen{\cdot, \cdot}_S$.
    Since $\gen{\cdot, \cdot}_S$ is an inner product, part (a) implies that $S$ is positive and invertible.
\end{solenum}

\begin{exercise}{7c24}
    Suppose $S$ and $T$ are positive operators on $V$.
    Prove that
    \[
        \nul(S + T) = \nul S \cap \nul T.
    \]
\end{exercise}

\begin{sol}
    \rightset Let $v \in \nul(S + T)$.
    Then $(S + T)v = 0$, so
    \[
        0 = \gen{(S + T)v, v} = \gen{Sv, v} + \gen{Tv, v}.
    \]
    Since $S$ and $T$ are positive, $\gen{Sv, v}$ and $\gen{Tv, v}$ are both nonnegative.
    Because their sum is zero, we must have $\gen{Sv, v} = 0$ and $\gen{Tv, v} = 0$.
    By 7.43, $Sv = 0$ and $Tv = 0$, so $v \in \nul S \cap \nul T$.
    
    \leftset Conversely, if $v \in \nul S \cap \nul T$, then $Sv = 0$ and $Tv = 0$, so $(S + T)v = Sv + Tv = 0$. Hence, $v \in \nul(S + T)$.
\end{sol}

\begin{exercise}{7c25}
    Let $T$ be the second derivative operator in \exref{7a31}(b).
    Show that $-T$ is a positive operator.
\end{exercise}

\begin{sol}
    Let $D \in \lin{V}$ by $Df = f'$.
    By \exref{7a31}(a), $D^\ast = -D$ for $Df = f'$.
    Since $Tf = f''$, we have $T = D^2$.
    Therefore, $-T = -D^2 = (-D)(D) = D^\ast D$, which is positive by 7.38(f).
\end{sol}